\section{Espacios vectoriales}
\label{sec:espacios_vectoriales}

\subsection{Definiciones b\'asicas}
\label{subsec:definiciones_ev}

\begin{definicion}[Espacio vectorial]
  \label{def:espacio_vectorial}
  Un \textbf{espacio vectorial lineal} sobre un campo $\F$ ($\R$ o $\C$) es un
  conjunto $V$ de objetos, denotados $\vb{a},\vb{b},\dots$ y llamados
  \emph{vectores}, tal que a cada par $\vb{a},\vb{b}\in V$ le corresponde un
  vector $\vb{a}+\vb{b}\in V$ que cumple las propiedades usuales de la suma
  (conmutatividad, asociatividad, neutro e inverso) y a cada
  $\al\in\F$, $\vb{a}\in V$ le corresponde $\al\vb{a}\in V$ compatible con
  ellas.
\end{definicion}

\begin{definicion}[Base]
  \label{def:base}
  Una \textbf{base} de $V$ es un conjunto $B$ de elementos linealmente
  independientes que generan $V$. Una base finita genera un espacio de
  dimensi\'on finita y una infinita un espacio de dimensi\'on infinita.
\end{definicion}

Si $B=\set{\vb{a}_i}_{i=1}^{n}$ es un conjunto de elementos linealmente
independientes, hablamos de un espacio de dimensi\'on $n$. La independencia
lineal significa que
\begin{equation}
  \label{eq:independencia_lineal}
  \sum_{i=1}^{n}\al_i\vh{e}_i=\vb{0}
  \qquad\iff\qquad
  \al_i=0\ \ \forall i .
\end{equation}

\begin{teorema}
  \label{teo:dimension}
  Todas las bases de un espacio de dimensi\'on finita tienen el mismo n\'umero
  de vectores linealmente independientes.
\end{teorema}

\begin{ejemplo}[Espacios finitos]
  \label{ej:espacios_finitos}
  $\R^2$, $\R^3$ y en general $\R^n$ son espacios vectoriales de dimensi\'on
  $2$, $3$ y $n$ respectivamente, con $n\in\Z^{+}\cup\set{0}$.
\end{ejemplo}

\begin{ejemplo}[Base infinita]
  \label{ej:base_infinita}
  Sea $V$ un espacio de funciones y $f\in V$ continua por pedazos y acotada. En
  un punto interior del dominio donde $f$ est\'a bien definida podemos
  aproximarla en una vecindad abierta de $x_0$ mediante su serie de Taylor,
  \begin{equation}
    \label{eq:taylor_1d}
    f(x)=\sum_{n=0}^{\infty}\frac{1}{n!}
      \eval{\dnv{f}{x}{n}}_{x_0}(x-x_0)^n ,
  \end{equation}
  es decir, por una aproximaci\'on polinomial. Para este curso una base de ese
  espacio de funciones es
  \begin{equation}
    \label{eq:base_polinomial}
    B=\set{x^n}_{n=0}^{\infty}.
  \end{equation}
\end{ejemplo}

\subsection{Representaci\'on geom\'etrica}
\label{subsec:representacion_geometrica}

Sean $\vb{a},\vb{b}\in V$; entonces $\vb{c}=\vb{a}+\vb{b}\in V$. Dada una base
$B=\set{\vh{e}_i}_{i=1}^{n}$ se tiene
\begin{equation}
  \label{eq:expansion_base}
  \vb{a}=\sum_{i=1}^{n}a_i\vh{e}_i,
  \qquad
  \vb{b}=\sum_{i=1}^{n}b_i\vh{e}_i,
\end{equation}
de modo que la suma y el producto por escalar act\'uan componente a componente,
\begin{equation}
  \label{eq:suma_componentes}
  \vb{a}+\vb{b}=\sum_{i=1}^{n}(a_i+b_i)\vh{e}_i,
  \qquad
  \al\vb{a}=\sum_{i=1}^{n}(\al a_i)\vh{e}_i ,
\end{equation}
y el vector se representa por su columna de coordenadas en esa base,
\begin{equation}
  \label{eq:columna_coordenadas}
  \vb{a}=
  \begin{bmatrix}
    a_1\\a_2\\\vdots\\a_n
  \end{bmatrix}_{e}.
\end{equation}

Para espacios de dimensi\'on infinita la idea es la misma. Si $f$ est\'a bien
definida en un dominio $D$ y admite expansi\'on de Taylor, la base
$B=\set{x^n}_{n=0}^{\infty}$ da \eqref{eq:taylor_1d}. Pero existen otras bases
$B=\set{g_n(x)}_{n=0}^{\infty}$ igualmente buenas, respecto de las cuales
\begin{equation}
  \label{eq:expansion_general}
  f(x)=\sum_{n=0}^{\infty}\hat{a}_n\,g_n(x).
\end{equation}
La Secci\'on~\ref{sec:fourier} explota esta libertad tomando exponenciales
complejas como base.

\subsection{Producto interno}
\label{subsec:producto_interno}

\begin{definicion}[Producto interno]
  \label{def:producto_interno}
  Un \textbf{producto interno} (o escalar) es una operaci\'on que asocia a dos
  elementos de $V$ un n\'umero del campo,
  \begin{equation}
    \label{eq:producto_interno}
    g:V\otimes V\longrightarrow\F,
    \qquad
    g(\vb{a},\vb{b})=\al\in\F .
  \end{equation}
\end{definicion}

\begin{ejemplo}[Representaci\'on en $\R^n$]
  \label{ej:pi_real}
  \begin{equation}
    \label{eq:pi_real}
    g(\vb{a},\vb{b})=\vb{a}\cdot\vb{b}=\vb{a}^\tp\vb{b}
      =a_1b_1+a_2b_2+\dots+a_nb_n=a_ib_i .
  \end{equation}
\end{ejemplo}

\begin{ejemplo}[Representaci\'on en $\C^n$]
  \label{ej:pi_complejo}
  \begin{equation}
    \label{eq:pi_complejo}
    g(\vb{a},\vb{b})=\vb{a}^\dg\vb{b}
      =\sqty{a_1,a_2,\dots,a_n}^{*}
      \begin{bmatrix}
        b_1\\b_2\\\vdots\\b_n
      \end{bmatrix},
  \end{equation}
  donde $*$ denota conjugaci\'on compleja.
\end{ejemplo}

\begin{observacion}[Par\'entesis complejo]
  Si $a\in\C$ lo representamos como $a=x+iy$ con $x,y\in\R$ e
  $i=\sqrt{-1}$, y definimos su \textbf{conjugado} $\bar{a}=a^{*}=x-iy$.
\end{observacion}

Las propiedades del producto interno complejo son:
\begin{enumerate}
  \item \textbf{Hermiticidad}: $g(\vb{a},\vb{b})=g(\vb{b},\vb{a})^{*}$.
  \item \textbf{Antilinealidad en la primera entrada}:
        \begin{equation}
          \label{eq:antilineal}
          g(\al\vb{a}+\be\vb{b},\vb{c})
            =\al^{*}g(\vb{a},\vb{c})+\be^{*}g(\vb{b},\vb{c}).
        \end{equation}
  \item \textbf{Linealidad en la segunda entrada}:
        \begin{equation}
          \label{eq:lineal}
          g(\vb{a},\de\vb{b}+\ga\vb{c})
            =\de\, g(\vb{a},\vb{b})+\ga\, g(\vb{a},\vb{c}).
        \end{equation}
\end{enumerate}

\begin{observacion}
  De la hermiticidad se sigue que $g(\vb{a},\vb{a})=g(\vb{a},\vb{a})^{*}$, es
  decir, $g(\vb{a},\vb{a})\in\R$ aun en el caso complejo. Adem\'as
  $g(\vb{a},\vb{a})\ge 0$, con igualdad si y s\'olo si $\vb{a}=\vb{0}$.
\end{observacion}

En el caso real la hermiticidad se reduce a simetr\'ia,
$g(\vb{a},\vb{b})=g(\vb{b},\vb{a})$, y la operaci\'on es lineal en ambas
entradas; por ejemplo
\begin{equation}
  \label{eq:bilineal_real}
  (\al\vb{a}+\be\vb{b})\cdot\vb{c}
    =\al\,\vb{a}\cdot\vb{c}+\be\,\vb{b}\cdot\vb{c}.
\end{equation}

\subsection{Norma}
\label{subsec:norma}

\begin{teorema}[Desigualdad de Schwarz]
  \label{teo:schwarz}
  Sean $\vb{a},\vb{b}\in V$ con un producto interno definido sobre $V$.
  Entonces
  \begin{equation}
    \label{eq:schwarz}
    \abs{g(\vb{a},\vb{b})}^2\le\norm{\vb{a}}^2\norm{\vb{b}}^2 .
  \end{equation}
\end{teorema}

\begin{definicion}[Norma]
  \label{def:norma}
  La \textbf{norma} de un elemento de $V$ con producto interno puede definirse a
  trav\'es del mismo producto interno, y cumple:
  \begin{itemize}
    \item $\norm{\vb{0}}=0$;
    \item $\norm{\vb{a}}\ge 0$, con $\norm{\vb{a}}=0\iff\vb{a}=\vb{0}$;
    \item $\norm{\al\vb{a}}=\abs{\al}\norm{\vb{a}}$ para todo $\al\in\C$;
    \item desigualdad del tri\'angulo,
          $\norm{\vb{a}+\vb{b}}\le\norm{\vb{a}}+\norm{\vb{b}}$.
  \end{itemize}
\end{definicion}

En particular se define la \textbf{norma $p$},
\begin{equation}
  \label{eq:norma_p}
  \norm{\vb{a}}_p=\qty{\sum_{i=1}^{n}\abs{a_i}^{p}}^{1/p},
\end{equation}
de la cual los casos relevantes son
\begin{equation}
  \label{eq:norma_2}
  \norm{\vb{a}}_2=\qty{\sum_{i=1}^{n}\abs{a_i}^{2}}^{1/2}
    =g(\vb{a},\vb{a})^{1/2}=(\vb{a}^\dg\vb{a})^{1/2}
\end{equation}
y
\begin{equation}
  \label{eq:norma_1}
  \norm{\vb{a}}_1=\sum_{i=1}^{n}\abs{a_i}.
\end{equation}

\subsection{Ortonormalizaci\'on de Gram--Schmidt}
\label{subsec:gram_schmidt}

Dado un conjunto linealmente independiente
$B=\set{\vb{v}_i}_{i=1}^{n}$ se puede construir a partir de \'el un conjunto
ortonormal. Se normaliza el primer vector,
\begin{equation}
  \label{eq:gs_primero}
  \vh{u}_1=\frac{\vb{v}_1}{\norm{\vb{v}_1}},
\end{equation}
se le resta a $\vb{v}_2$ su proyecci\'on sobre $\vh{u}_1$,
\begin{equation}
  \label{eq:gs_proyeccion}
  \tilde{\vb{v}}_2=\vb{v}_2-\qty{\vb{v}_2\cdot\vh{u}_1}\vh{u}_1,
\end{equation}
y se normaliza el resultado,
\begin{equation}
  \label{eq:gs_segundo}
  \vh{u}_2=\frac{\tilde{\vb{v}}_2}{\norm{\tilde{\vb{v}}_2}} .
\end{equation}
En general, habiendo construido $\vh{u}_1,\dots,\vh{u}_{k-1}$,
\begin{equation}
  \label{eq:gs_general}
  \tilde{\vb{v}}_k=\vb{v}_k-\sum_{j=1}^{k-1}
    \qty{\vb{v}_k\cdot\vh{u}_j}\vh{u}_j,
  \qquad
  \vh{u}_k=\frac{\tilde{\vb{v}}_k}{\norm{\tilde{\vb{v}}_k}} .
\end{equation}

\subsection{Producto cu\~na en $\R^3$}
\label{subsec:producto_cuna_3d}

\begin{definicion}[Producto cu\~na]
  \label{def:producto_cuna}
  Sean $\vb{a},\vb{b}\in V\subset\R^3$. Definimos la operaci\'on
  $\wedge:V\otimes V\rightarrow V$ tal que
  $\vb{a}\wedge\vb{b}=\vb{c}$, donde $\vb{c}$ es perpendicular a ambos
  vectores.
\end{definicion}

Sus propiedades son:
\begin{enumerate}
  \item $\vb{a}\wedge\vb{b}=\vb{0}$ si $\vb{b}\parallel\vb{a}$ o
        $\vb{b}=\vb{0}$;
  \item $\vb{a}\wedge\vb{b}=-\vb{b}\wedge\vb{a}$ (antisimetr\'ia);
  \item $\vb{a}\wedge(\al\vb{b}+\be\vb{c})
         =\al\,\vb{a}\wedge\vb{b}+\be\,\vb{a}\wedge\vb{c}$;
  \item $(\al\vb{a}+\be\vb{b})\wedge\vb{c}
         =\al\,\vb{a}\wedge\vb{c}+\be\,\vb{b}\wedge\vb{c}$;
  \item $(\al\vb{a})\wedge\vb{b}=\al\qty{\vb{a}\wedge\vb{b}}$.
\end{enumerate}
El sentido de $\vb{c}$ lo fija la regla de la mano derecha.

\subsubsection{Interpretaci\'on f\'isica}
\label{subsubsec:torque}

El momento de una fuerza es
\begin{equation}
  \label{eq:torque}
  \boldsymbol{\ta}=\vb{r}\wedge\vb{F},
\end{equation}
con magnitud
\begin{equation}
  \label{eq:magnitud_torque}
  \norm{\boldsymbol{\ta}}=\norm{\vb{r}\wedge\vb{F}}
    =r F_{\perp}=\norm{\vb{r}}\norm{\vb{F}}\sin\theta .
\end{equation}

En general
\begin{equation}
  \label{eq:norma_cuna}
  \norm{\vb{a}\wedge\vb{b}}=\norm{\vb{a}}\norm{\vb{b}}\sin\theta,
\end{equation}
que es dos veces el \'area del tri\'angulo generado por $\vb{a}$ y $\vb{b}$, o
equivalentemente el \'area del paralelogramo que determinan. Por eso
$\vb{c}=\vb{a}\wedge\vb{b}=\vb{A}$ se interpreta como el \textbf{vector de
\'area} de esa caja bidimensional: tiene magnitud igual al \'area y direcci\'on
normal a ella.

\begin{observacion}
  Toda superficie tiene entonces un campo de vectores de \'area, uno sobre cada
  parche, y ese campo dota al espacio de una orientaci\'on. Como consecuencia,
  para una curva simple cerrada se dice que se recorre en \textbf{sentido
  positivo} si el interior de la curva queda del lado izquierdo.
\end{observacion}
